\section{Additional results on irreducibility}

\begin{theorem}[Orthogonal projections and involutive-algebra projections]
    \label{thm:orthogonal_projection_star_projection_channel}
    \lean{IsOrthogonalProjection.isStarProjection}
    \lean{IsStarProjection.isOrthogonalProjection}
    \leanok
    \uses{def:orthogonal_projection_channel}
    For a complex square matrix, the two descriptions
    $P = P^\dagger$, $P^2=P$ and $P^* = P$, $P^2=P$ are equivalent.
\end{theorem}

\begin{proof}\leanok
    Since the involution on $\MN{D}$ is conjugate transpose,
    $P^* = P^\dagger$. Therefore
    \[
        P^\dagger = P,\; P^2 = P
        \iff
        P^* = P,\; P^2 = P.
    \]
\end{proof}

\begin{theorem}[Coisometric transport of a star projection]
    \label{thm:star_projection_coisometric_transport}
    \lean{IsStarProjection.conjTranspose_mul_mul_of_mul_conjTranspose_eq_one}
    \leanok
    Let $Q$ be a self-adjoint idempotent on a finite-dimensional space, and
    let $U$ be a coisometry onto that space, so that $UU^\dagger=\Id$.  Then
    $U^\dagger Q U$ is self-adjoint and idempotent.
\end{theorem}

\begin{proof}\leanok
    Self-adjointness follows from
    $(U^\dagger Q U)^\dagger=U^\dagger Q^\dagger U=U^\dagger Q U$.  Moreover,
    \[
        (U^\dagger Q U)^2
        =U^\dagger Q(UU^\dagger)QU
        =U^\dagger Q^2U
        =U^\dagger Q U.
    \]
\end{proof}

\begin{lemma}[Pairwise orthogonality from a projection resolution]
    \label{lem:orthogonal_projection_mul_eq_zero_of_sum_eq_one}
    \lean{orthogonalProjection_mul_eq_zero_of_sum_eq_one}
    \leanok
    \uses{def:orthogonal_projection_channel}
    Let $(Q_k)_k$ be a finite family of orthogonal projections satisfying
    $\sum_k Q_k=\Id$.  If $k\ne\ell$, then $Q_kQ_\ell=0$.
\end{lemma}

\begin{proof}\leanok
    Multiplying the resolution of the identity on both sides by $Q_k$ gives
    \[
        \sum_j Q_kQ_jQ_k=Q_k.
    \]
    Each summand is positive semidefinite because
    $Q_kQ_jQ_k=(Q_jQ_k)^*(Q_jQ_k)$, while the $j=k$ summand is already $Q_k$.
    Subtracting that term and taking traces gives
    \[
        0=\sum_{j\ne k}\tr\bigl((Q_jQ_k)^*(Q_jQ_k)\bigr)
         =\sum_{j\ne k}\lVert Q_jQ_k\rVert_{\mathrm F}^{2}.
    \]
    Hence $Q_jQ_k=0$ for $j\ne k$, and taking adjoints gives $Q_kQ_j=0$.
\end{proof}

\begin{theorem}[Complement of an orthogonal projection]
    \label{thm:orthogonal_projection_complement_channel}
    \lean{IsOrthogonalProjection.one_sub}
    \leanok
    \uses{def:orthogonal_projection_channel,
        thm:orthogonal_projection_star_projection_channel}
    If $P \in \MN{D}$ is an orthogonal projection, then $\Id-P$ is an
    orthogonal projection.
\end{theorem}

\begin{proof}\leanok
    In the involutive algebra $\MN{D}$, $P$ satisfies $P^\dagger=P$ and
    $P^2=P$. Hence
    \[
        (\Id-P)^\dagger=\Id-P,\qquad
        (\Id-P)^2=\Id-2P+P^2=\Id-P .
    \]
    Thus $\Id-P$ is an orthogonal projection.
\end{proof}

\begin{theorem}[Injectivity implies irreducibility]\label{thm:injective_irreducible}
    \lean{MPSTensor.injective_implies_irreducibleCP}
    \leanok
    \uses{def:injective, def:transfer_map, def:irreducible_cp}
    If $A$ is an injective MPS tensor, then its transfer map $\E_A$
    is irreducible.
\end{theorem}

\begin{proof}
    \leanok
    If $P$ is an invariant projection for $\E_A$,
    then $(\Id - P) A^i P = 0$ for all~$i$.
    Since the $\{A^i\}$ span $\MN{D}$,
    $(\Id - P) M P = 0$ for all $M$,
    so $P = 0$ or $P = \Id$.
\end{proof}

\begin{theorem}[Transfer map is CP]\label{thm:transfer_cp}
    \lean{MPSTensor.transferMap_isCPMap}
    \leanok
    \uses{def:transfer_map, def:cp_map}
    The transfer map $\E_A(X) = \sum_i A^i X (A^i)^\dagger$
    is completely positive.
\end{theorem}

\begin{proof}\leanok
    It is already written in Kraus form, with Kraus operators $\{A^i\}$.
\end{proof}

\begin{theorem}[Transfer map is a channel]\label{thm:transfer_channel}
    \lean{MPSTensor.transferMap_isChannel}
    \leanok
    \uses{def:transfer_map, def:channel}
    If $\sum_i (A^i)^\dagger A^i = \Id$, then the transfer map
    $\E_A(X) = \sum_i A^i X (A^i)^\dagger$
    is a quantum channel (CPTP).
\end{theorem}

\begin{proof}\leanok\uses{thm:transfer_cp}
    Complete positivity is the observation of Theorem~\ref{thm:transfer_cp}.
    Trace preservation:
    $\tr(\E_A(X)) = \sum_i \tr(A^i X (A^i)^\dagger)
        = \tr((\sum_i (A^i)^\dagger A^i) X) = \tr(X)$
    by cyclicity and the normalization hypothesis.
\end{proof}

\section{Ces\`aro fixed-point theorem}

\begin{definition}[Ces\`aro mean]\label{def:cesaro_mean}
    \lean{cesaroMean}
    \leanok
    The \emph{Ces\`aro mean} of a linear map $E$ at order $N$,
    applied to a starting matrix $\rho_0$, is
    \[
        \sigma_N =   \frac{1}{N} \sum_{n=0}^{N-1} E^n(\rho_0).
    \]
\end{definition}

\begin{theorem}[Ces\`aro telescope]\label{thm:cesaro_telescope}
    \lean{cesaroMean_telescope}
    \leanok
    \uses{def:cesaro_mean}
    $E(\sigma_N) - \sigma_N = \frac{1}{N}(E^N(\rho_0) - \rho_0)$.
\end{theorem}

\begin{proof}\leanok\uses{def:cesaro_mean}
    Telescoping: $E(\sigma_N) = \frac{1}{N}\sum_{n=1}^{N} E^n(\rho_0)$
    and subtracting $\sigma_N$ cancels all but the first and last terms.
\end{proof}

\begin{definition}[Bounded-orbit endomorphism]
    \label{def:has_bounded_orbits}
    \lean{LinearMap.HasBoundedOrbits}
    \leanok
    Let $\mathbb K$ denote either $\R$ or $\C$, and let $V$ be a normed
    vector space over $\mathbb K$. A linear endomorphism $f:V\to V$ has
    \emph{bounded orbits} if, for every $x\in V$, the set
    \[
        \{f^n x:n\in\N\}
    \]
    is bounded.
\end{definition}

\begin{lemma}[Ces\`aro averages of range elements]
    \label{lem:mean_ergodic_range_limit}
    \lean{LinearMap.HasBoundedOrbits.tendsto_birkhoffAverage_of_mem_range_sub_one}
    \leanok
    \uses{def:has_bounded_orbits}
    Let $f:V\to V$ have bounded orbits. If
    $x\in\range(f-\Id_V)$, then
    \[
        A_N(x)\longrightarrow 0.
    \]
\end{lemma}

\begin{proof}\leanok
    \uses{def:has_bounded_orbits}
    Write $x=(f-\Id_V)y$. Telescoping gives
    \[
        A_N(x)=\frac{f^N y-y}{N}.
    \]
    Boundedness of the orbit of $y$ makes the numerator bounded, so the
    right-hand side tends to zero.
\end{proof}

\begin{lemma}[Complementarity of the fixed-point space]
    \label{lem:mean_ergodic_complementarity}
    \lean{LinearMap.HasBoundedOrbits.disjoint_ker_sub_one_range_sub_one,
        LinearMap.HasBoundedOrbits.isCompl_ker_sub_one_range_sub_one}
    \leanok
    \uses{def:has_bounded_orbits, lem:mean_ergodic_range_limit}
    Let $V$ be finite-dimensional and let $f:V\to V$ have bounded orbits.
    Then
    \[
        V=\ker(f-\Id_V)\oplus\range(f-\Id_V).
    \]
\end{lemma}

\begin{proof}\leanok
    \uses{def:has_bounded_orbits, lem:mean_ergodic_range_limit}
    If $z$ belongs to both summands, then $A_N(z)=z$ because
    $z\in\ker(f-\Id_V)$, while Lemma~\ref{lem:mean_ergodic_range_limit} gives
    $A_N(z)\to 0$. Hence $z=0$. Rank--nullity then gives the asserted
    complementarity.
\end{proof}

\begin{definition}[Mean-ergodic projection]
    \label{def:mean_ergodic_projection}
    \lean{LinearMap.meanErgodicProjection}
    \leanok
    \uses{def:has_bounded_orbits, lem:mean_ergodic_complementarity}
    Let $V$ be finite-dimensional and let $f:V\to V$ have bounded orbits.
    The \emph{mean-ergodic projection} $P_f:V\to V$ is the projection onto
    $\ker(f-\Id_V)$ along $\range(f-\Id_V)$.
\end{definition}

\begin{theorem}[Finite-dimensional mean-ergodic theorem]
    \label{thm:mean_ergodic_bounded_orbits}
    \lean{LinearMap.HasBoundedOrbits.tendsto_birkhoffAverage_meanErgodicProjection,
        LinearMap.HasBoundedOrbits.range_meanErgodicProjection,
        LinearMap.HasBoundedOrbits.meanErgodicProjection_apply_eq_self_iff,
        LinearMap.HasBoundedOrbits.isIdempotentElem_meanErgodicProjection,
        LinearMap.HasBoundedOrbits.meanErgodicProjection_apply_meanErgodicProjection,
        LinearMap.HasBoundedOrbits.comp_meanErgodicProjection,
        LinearMap.HasBoundedOrbits.meanErgodicProjection_comp}
    \leanok
    \uses{def:has_bounded_orbits, def:mean_ergodic_projection}
    Let $V$ be finite-dimensional and let $f:V\to V$ have bounded orbits.
    For every $x\in V$, the Ces\`aro averages
    \[
        A_N(x)=\frac{1}{N}\sum_{n=0}^{N-1}f^n x
    \]
    converge to $P_f x$. Moreover,
    \[
        \range(P_f)=\ker(f-\Id_V),\qquad
        P_f^2=P_f,\qquad fP_f=P_f=P_f f,
    \]
    and $P_f x=x$ if and only if $f x=x$.
    The complex matrix specialization underlying
    \cite[Equation~(6.14)]{Wolf2012Quantum} is given below.
\end{theorem}

\begin{proof}\leanok
    \uses{def:has_bounded_orbits, lem:mean_ergodic_range_limit,
        lem:mean_ergodic_complementarity, def:mean_ergodic_projection}
    By Lemma~\ref{lem:mean_ergodic_complementarity}, decompose
    \[
        x=P_f x+(x-P_f x),\qquad
        P_f x\in\ker(f-\Id_V),\qquad
        x-P_f x\in\range(f-\Id_V).
    \]
    The fixed component has constant Ces\`aro averages, whereas the averages of
    the second component tend to zero by
    Lemma~\ref{lem:mean_ergodic_range_limit}. Thus $A_N(x)\to P_f x$. Since
    \[
        P_f x\in\ker(f-\Id_V),
        \qquad f(P_f x)=P_f x,
    \]
    we have $fP_f=P_f$. Moreover,
    \[
        (f-\Id_V)x\in\range(f-\Id_V),
        \qquad P_f\bigl((f-\Id_V)x\bigr)=0,
    \]
    so $P_f(fx)=P_f x$ and hence $P_f f=P_f$. The range and idempotence
    identities follow directly from the projection construction.
\end{proof}

\begin{theorem}[Positive mean-ergodic projection]
    \label{thm:positive_mean_ergodic_projection}
    \lean{IsPositiveMap.meanErgodicProjection_isPositiveMap}
    \leanok
    \uses{def:positive_map, def:mean_ergodic_projection,
        thm:mean_ergodic_bounded_orbits}
    Let $T:\MN{D}\to\MN{D}$ be positive and have bounded orbits. Then its
    mean-ergodic projection $P_T$ is positive.
\end{theorem}

\begin{proof}\leanok
    Every Ces\`aro average is positive. Therefore, for $X\ge 0$, closedness of
    the positive cone and convergence of the averages give $P_T(X)\ge 0$.
\end{proof}

\begin{theorem}[Trace-preserving mean-ergodic projection]
    \label{thm:trace_preserving_mean_ergodic_projection}
    \lean{IsTracePreservingMap.meanErgodicProjection_isTracePreservingMap}
    \leanok
    \uses{def:trace_preserving, def:mean_ergodic_projection,
        thm:mean_ergodic_bounded_orbits}
    Let $T:\MN{D}\to\MN{D}$ be trace-preserving and have bounded orbits.
    Then its mean-ergodic projection $P_T$ is trace-preserving.
\end{theorem}

\begin{proof}\leanok
    Every nonempty Ces\`aro average preserves the trace. Convergence of the
    averages and continuity of the trace therefore give
    $\tr(P_T(X))=\tr(X)$.
\end{proof}

\begin{theorem}[Identity under the mean-ergodic projection]
    \label{thm:unital_mean_ergodic_projection}
    \lean{LinearMap.HasBoundedOrbits.meanErgodicProjection_one}
    \leanok
    \uses{def:mean_ergodic_projection, thm:mean_ergodic_bounded_orbits}
    Let $T:\MN{D}\to\MN{D}$ have bounded orbits. If $T(\Id)=\Id$, then
    $P_T(\Id)=\Id$.
\end{theorem}

\begin{proof}\leanok
    The identity is a fixed point of $T$, so the fixed-point characterization
    of the mean-ergodic projection gives the conclusion.
\end{proof}

\begin{theorem}[Bounded orbits from trace nonincrease]
    \label{thm:positive_trace_nonincreasing_bounded_orbits}
    \lean{IsPositiveMap.hasBoundedOrbits_of_traceNonincreasing}
    \leanok
    \uses{def:positive_map, def:has_bounded_orbits}
    Let $T:\MN{D}\to\MN{D}$ be positive and trace nonincreasing.  Then the
    forward orbit $\{T^n(X):n\geq 0\}$ is bounded for every $X\in\MN{D}$.
    This is the trace-nonincreasing form of
    \cite[Proposition~6.3]{Wolf2012Quantum}.
\end{theorem}

\begin{proof}\leanok
    For $X\succeq 0$, positivity and induction on $n$ give
    \[
        T^n(X)\succeq 0,
        \qquad
        \tr(T^n(X))\leq\tr(X).
    \]
    Thus the orbit lies in a bounded trace section of the positive cone.
    Apply this to the positive and negative parts of a Hermitian matrix, and
    then write an arbitrary matrix as a complex linear combination of two
    Hermitian matrices.
\end{proof}

\begin{theorem}[Positive trace-preserving mean-ergodic projection]
    \label{thm:positive_trace_preserving_mean_ergodic_projection}
    \lean{IsPositiveMap.hasBoundedOrbits_of_tracePreserving,
        IsPositiveMap.range_meanErgodicProjection_of_tracePreserving,
        IsPositiveMap.meanErgodicProjection_apply_eq_self_iff_of_tracePreserving}
    \leanok
    \uses{def:positive_map, def:trace_preserving,
        def:mean_ergodic_projection,
        thm:positive_mean_ergodic_projection,
        thm:trace_preserving_mean_ergodic_projection,
        thm:unital_mean_ergodic_projection}
    Let $T:\MN{D}\to\MN{D}$ be positive and trace-preserving. Then $T$ has
    bounded orbits, and its mean-ergodic projection $P_T$ is positive and
    trace-preserving. Its range is precisely the fixed-point space of $T$:
    \[
        \range(P_T)=\ker(T-\Id),
        \qquad
        P_T(X)=X \iff T(X)=X.
    \]
    If $T(\Id)=\Id$, then $P_T(\Id)=\Id$.
    This is the Ces\`aro projection $T_\infty$ of
    \cite[Proposition~6.3 and Equation~(6.14)]{Wolf2012Quantum}.
\end{theorem}

\begin{proof}\leanok
    \uses{thm:mean_ergodic_bounded_orbits,
        thm:positive_trace_nonincreasing_bounded_orbits,
        thm:positive_mean_ergodic_projection,
        thm:trace_preserving_mean_ergodic_projection,
        thm:unital_mean_ergodic_projection}
    Trace preservation implies trace nonincrease, so
    Theorem~\ref{thm:positive_trace_nonincreasing_bounded_orbits} gives
    bounded orbits.  Theorem~\ref{thm:mean_ergodic_bounded_orbits} then gives
    the projection and its fixed-point characterization.  Its positivity,
    trace preservation, and behavior on the identity follow from the three
    preceding theorems.
\end{proof}

\begin{theorem}[Trace preservation and the trace adjoint]
    \label{thm:trace_preserving_iff_trace_adjoint_unital}
    \lean{isTracePreservingMap_iff_traceAdjointMap_one}
    \leanok
    \uses{def:trace_preserving, def:trace_pairing_adjoint}
    On any finite full matrix algebra, a linear endomorphism $S$ is
    trace-preserving if and only if its trace-pairing adjoint fixes the
    identity:
    \[
        (\forall X,\; \tr(S(X))=\tr(X))
        \quad\Longleftrightarrow\quad
        S^*(\Id)=\Id.
    \]
\end{theorem}

\begin{proof}\leanok
    \uses{thm:trace_pairing_adjoint_identity, thm:trace_nondeg}
    The trace-pairing identity gives
    \[
        \tr(S^*(\Id)X)=\tr(S(X)).
    \]
    The equivalence follows from nondegeneracy of the trace pairing.
\end{proof}

\begin{theorem}[Adjoint mean-ergodic projection]
    \label{thm:adjoint_mean_ergodic_projection}
    \lean{LinearMap.HasBoundedOrbits.traceAdjointMap_meanErgodicProjection_apply_eq_self_iff,
        LinearMap.HasBoundedOrbits.traceAdjointMap_meanErgodicProjection_apply_self,
        LinearMap.HasBoundedOrbits.range_traceAdjointMap_meanErgodicProjection}
    \leanok
    \uses{def:mean_ergodic_projection, thm:mean_ergodic_bounded_orbits}
    Let $T:\MN{D}\to\MN{D}$ have bounded orbits, let $P_T$ be its
    mean-ergodic projection, and let $P_T^*$ be the trace-pairing adjoint of
    $P_T$. Then
    \[
        (P_T^*)^2=P_T^*,
        \qquad
        \range(P_T^*)=\ker(T^*-\Id),
    \]
    and
    \[
        P_T^*(Y)=Y\quad\Longleftrightarrow\quad T^*(Y)=Y.
    \]
\end{theorem}

\begin{proof}\leanok
    \uses{thm:trace_pairing_adjoint_identity, thm:trace_nondeg,
        thm:mean_ergodic_bounded_orbits}
    Since $P_T^2=P_T$, the trace-pairing identity gives
    $(P_T^*)^2=P_T^*$.  If $T^*(Y)=Y$, decompose
    \[
        X-P_T(X)=(T-\Id)(Z).
    \]
    Then
    \[
        \tr\bigl(Y(X-P_T(X))\bigr)
        =\tr\bigl((T^*(Y)-Y)Z\bigr)=0.
    \]
    Nondegeneracy of the trace pairing gives $P_T^*(Y)=Y$.  Conversely,
    $P_TT=P_T$ implies $T^*P_T^*=P_T^*$, so every fixed point of $P_T^*$ is a
    fixed point of $T^*$:
    \[
        P_T^*(Y)=Y
        \quad\Longrightarrow\quad
        T^*(Y)=T^*(P_T^*(Y))=P_T^*(Y)=Y.
    \]
    Hence
    \[
        Y\in\range(P_T^*)
        \quad\Longleftrightarrow\quad P_T^*(Y)=Y
        \quad\Longleftrightarrow\quad T^*(Y)=Y
        \quad\Longleftrightarrow\quad Y\in\ker(T^*-\Id).
    \]
    The first equivalence uses idempotence, and the second is the fixed-point
    equivalence proved above.
\end{proof}

\begin{theorem}[Positive unital adjoint Ces\`aro retraction]
    \label{thm:positive_unital_adjoint_mean_ergodic_retraction}
    \lean{IsPositiveMap.traceAdjoint_meanErgodicProjection_isPositiveUnitalRetraction}
    \leanok
    \uses{thm:positive_trace_preserving_mean_ergodic_projection,
        thm:trace_preserving_iff_trace_adjoint_unital,
        thm:adjoint_mean_ergodic_projection,
        thm:trace_pairing_adjoint_positive}
    Let $T:\MN{D}\to\MN{D}$ be positive and trace-preserving.  The trace
    adjoint $P_T^*$ of its mean-ergodic projection is positive, unital, and
    idempotent.  It is a retraction onto the adjoint fixed-point space:
    \[
        P_T^*(\Id)=\Id,
        \qquad (P_T^*)^2=P_T^*,
        \qquad \range(P_T^*)=\ker(T^*-\Id),
    \]
    and $P_T^*(Y)=Y$ if and only if $T^*(Y)=Y$.
    This is the adjoint projection used in the proof of
    \cite[Theorem~6.14]{Wolf2012Quantum}.
\end{theorem}

\begin{proof}\leanok
    Theorem~\ref{thm:positive_trace_preserving_mean_ergodic_projection} makes
    $P_T$ positive and trace-preserving.  Positivity passes to its trace
    adjoint, while
    Theorem~\ref{thm:trace_preserving_iff_trace_adjoint_unital} gives
    $P_T^*(\Id)=\Id$.  The remaining assertions are
    Theorem~\ref{thm:adjoint_mean_ergodic_projection}.
\end{proof}

\begin{theorem}[Hermitian fixed-point decomposition]\label{thm:hermitian_fixedpoint_parts}
    \lean{IsPositiveMap.posPart_negPart_fixed_of_hermitian_fixedPoint}
    \leanok
    \uses{def:positive_map, def:trace_preserving}
    Let $E$ be a positive trace-preserving linear map and let $X = E(X)$ be a Hermitian
    fixed point. Let $P_+=X^+$ and $P_-=X^-$ be its canonical positive and negative
    parts. Thus $X=P_+-P_-$, $P_\pm\geq0$, and $P_+P_-=P_-P_+=0$.
    Then $E(P_+) = P_+$ and $E(P_-) = P_-$.
    This is \cite[Proposition~6.8]{Wolf2012Quantum}.
\end{theorem}

\begin{proof}
    \leanok
    The fixed-point equation gives a common difference
    \[
        Y=E(P_+)-P_+=E(P_-)-P_-.
    \]
    Hence $P_++Y$ and $P_-+Y$ are positive semidefinite, while trace
    preservation gives $\tr(Y)=0$. Diagonalize $P_+$. Since
    $P_+P_-=P_-P_+=0$, each eigenvector belongs to the kernel of $P_+$ or to
    the kernel of $P_-$. Positivity of the corresponding matrix $P_++Y$ or
    $P_-+Y$ shows that every diagonal coefficient of $Y$ in this basis is
    nonnegative. Their sum is $\tr(Y)=0$, so every such coefficient vanishes.
    A positive semidefinite matrix with a zero diagonal coefficient annihilates
    the corresponding basis vector. Applying this to $P_++Y$ or $P_-+Y$ shows
    that $Y$ annihilates the entire eigenbasis. Thus $Y=0$, and consequently
    $E(P_\pm)=P_\pm$.
\end{proof}

\begin{theorem}[Ces\`aro fixed-point theorem]\label{thm:cesaro_fixedpoint}
    \lean{IsChannel.exists_posSemidef_fixedPoint}
    \leanok
    \uses{def:channel, def:density_matrices, def:cesaro_mean}
    Every quantum channel $E : \MN{D} \to \MN{D}$ (with $D \ge 1$)
    has a nonzero positive semidefinite fixed point:
    there exists $\rho \ge 0$, $\rho \neq 0$, with $E(\rho) = \rho$.
\end{theorem}

\begin{proof}
    \leanok
    \uses{thm:density_compact, thm:density_nonempty,
        thm:channel_preserves_density, thm:cesaro_telescope}
    Start with any $\rho_0 \in \mc{D}_D$
    (Theorem~\ref{thm:density_nonempty}).
    Each Ces\`aro mean $\sigma_N$ lies in $\mc{D}_D$
    by applying channel preservation to every iterate, adding positive
    semidefinite matrices, rescaling by the positive average factor, and using
    trace linearity.
    By compactness (Theorem~\ref{thm:density_compact}),
    extract a convergent subsequence
    $\sigma_{N_k} \to \rho$.
    The telescope identity (Theorem~\ref{thm:cesaro_telescope})
    gives
    $\|E(\sigma_N) - \sigma_N\| = \frac{1}{N}\|E^N(\rho_0) - \rho_0\| \to 0$,
    so $E(\rho) = \rho$ by continuity.
    In~\cite[Theorem~6.11]{Wolf2012Quantum}, existence is proved
    via Brouwer's fixed-point theorem; we use the Ces\`aro
    argument instead.
\end{proof}

\section{Fixed-point projection}

\begin{definition}[Fixed-point projection]\label{def:fixed_point_proj}
    \lean{fixedPointProj}
    \leanok
    Given a linear map $E : \MN{D} \to \MN{D}$ and a fixed point
    $\rho$ with $E(\rho) = \rho$ and $\tr(\rho) \neq 0$,
    the \emph{fixed-point projection} is the rank-one map
    \[
        P(X) =  \frac{\tr(X)}{\tr(\rho)} \rho.
    \]
    This projects onto the span of $\rho$ along the kernel of the
    trace functional. When the fixed point is unique up to scaling,
    this span is the full fixed-point space.
\end{definition}

\begin{theorem}[Power decomposition]\label{thm:pow_decomp}
    \lean{pow_eq_fixedPointProj_add_compl_pow}
    \leanok
    \uses{def:fixed_point_proj, def:trace_preserving}
    If $E$ is trace-preserving and $E(\rho) = \rho$, then
    for $n \ge 1$,
    \[
        E^n =  P + (E - P)^n,
    \]
    where $P$ is the fixed-point projection.
\end{theorem}

\begin{proof}
    \leanok
    \uses{def:fixed_point_proj}
    $P$ is idempotent ($P^2 = P$), the fixed-point relation
    $E(\rho) = \rho$ gives $E \circ P = P$, and
    trace preservation gives $P \circ E = P$.
    These imply $(E-P) \circ P = 0$ and $P \circ (E-P) = 0$,
    so in the binomial expansion of $(P + (E-P))^n$
    all cross terms vanish, leaving $E^n = P^n + (E-P)^n = P + (E-P)^n$.
\end{proof}

\section{Additional results on peripheral spectrum and primitivity}

\begin{theorem}[$1$ is a peripheral eigenvalue]\label{thm:one_peripheral}
    \lean{one_mem_peripheralEigenvalues}
    \leanok
    \uses{def:peripheral_eigenvalues}
    If $E$ has a nonzero fixed point $\rho \neq 0$ with $E(\rho) = \rho$,
    then $1$ is a peripheral eigenvalue of $E$.
\end{theorem}

\begin{proof}\leanok
    $\rho$ is an eigenvector with eigenvalue~$1$, and $|1| = 1$.
\end{proof}

\begin{theorem}[Finiteness of peripheral eigenvalues]\label{thm:peripheral_finite}
    \lean{peripheralEigenvalues_finite}
    \leanok
    \uses{def:peripheral_eigenvalues}
    On a finite-dimensional space, the set of peripheral eigenvalues
    is finite.
\end{theorem}

\begin{proof}\leanok
    Peripheral eigenvalues are a subset of all eigenvalues,
    which is a finite set in finite dimensions.
\end{proof}

\begin{theorem}[Power-stable peripheral eigenvalues are roots of unity]\label{thm:peripheral_roots_of_unity}
    \lean{peripheral_isRootOfUnity_of_pow_eigenvalue}
    \leanok
    \uses{def:peripheral_eigenvalues}
    Let $E$ be a linear endomorphism on a finite-dimensional
    space.
    If $\mu$ is a peripheral eigenvalue of $E$
    and $\mu^n$ is an eigenvalue of $E$ for every $n \ge 1$,
    then $\mu$ is a root of unity.
\end{theorem}

\begin{proof}
    \leanok
    Since $\mu^n$ is an eigenvalue for every positive~$n$,
    the pigeonhole principle on the finite eigenvalue set gives
    $\mu^a = \mu^b$ for some $a \neq b$, hence $\mu^{|a-b|} = 1$.
\end{proof}

\begin{theorem}[Peripheral powers imply root of unity]\label{thm:peripheral_powers}
    \lean{peripheral_isRootOfUnity_of_closed_powers}
    \leanok
    \uses{def:peripheral_eigenvalues}
    If the set of peripheral eigenvalues of a linear
    endomorphism~$E$ on a finite-dimensional space is
    \emph{closed under powers}
    ($\mu \in \mathrm{peripheral}(E)$
    and $n \ge 1$ imply
    $\mu^n \in \mathrm{peripheral}(E)$),
    then every peripheral eigenvalue is a root of unity.
\end{theorem}

\begin{proof}
    \leanok
    \uses{thm:peripheral_roots_of_unity}
    The closure hypothesis provides
    $\mu^n \in \mathrm{peripheral}(E)$ for all positive~$n$,
    which in particular means $\mu^n$ is an eigenvalue.
    Theorem~\ref{thm:peripheral_roots_of_unity} then gives
    $\mu^p = 1$ for some $p > 0$.
\end{proof}

\subsection{Periodicity removal by powering}\label{subsec:periodicity_removal_power}

\begin{remark}\label{rem:periodicity_removal_spectral}\leanok
    The results in this subsection are purely spectral: they use only
    finiteness of a set of complex numbers and the spectral mapping theorem
    for powers. In particular, they do not rely on complete positivity.
    For transfer maps, replacing a tensor by its $p$-blocked tensor replaces
    $\E_A$ by $\E_A^p$, so this powering step is the operator-theoretic
    form of blocking.
\end{remark}

\begin{lemma}[Common exponent for a finite set of roots of unity]\label{lem:common_power_eq_one}
    \lean{exists_common_power_eq_one_of_finite}
    \leanok
    Let $s \subseteq \C$ be a finite set.
    Assume that for every $\mu \in s$ there exists an exponent $p_\mu > 0$
    with $\mu^{p_\mu} = 1$.
    Then there exists $p > 0$ such that $\mu^p = 1$ for all $\mu \in s$.
\end{lemma}

\begin{proof}\leanok
    Take $p$ to be the least common multiple of the exponents $p_\mu$.
\end{proof}

\begin{lemma}[Powering collapses peripheral eigenvalues]\label{lem:peripheral_pow_singleton}
    \lean{peripheralEigenvalues_pow_eq_singleton}
    \leanok
    \uses{def:peripheral_eigenvalues}
    Let $E : \MN{D} \to \MN{D}$ be a linear map, and assume that $E$ has a
    nonzero fixed point $\rho \neq 0$.
    Let $p > 0$.
    If every peripheral eigenvalue $\mu$ of $E$ satisfies $\mu^p = 1$,
    then $\mathrm{peripheral}(E^p) = \{1\}$.
\end{lemma}

\begin{proof}\leanok
    \uses{thm:one_peripheral}
    The fixed point $\rho$ gives $1 \in \mathrm{peripheral}(E^p)$.
    For the reverse inclusion, let $\nu \in \mathrm{peripheral}(E^p)$.
    Spectral mapping gives $\nu = \mu^p$ for some $\mu \in \sigma(E)$.
    From $|\nu| = 1$ and $p > 0$ we obtain $|\mu| = 1$, hence
    $\mu \in \mathrm{peripheral}(E)$.
    The hypothesis $\mu^p = 1$ then forces $\nu = 1$.
\end{proof}

For a normalized MPS tensor, the transfer map is naturally trace-preserving but need
not be unital. In general one should not expect a single gauge choice to make it
both unital and trace-preserving.

\begin{remark}[Bi-canonical special case]\label{rem:bicanonical_special_case}\leanok
    When a Kraus map happens to be both unital and trace-preserving
    (the ``bi-canonical'' case), the Kadison--Schwarz equality at peripheral
    eigenvectors (Theorem~\ref{thm:ks_peripheral_channels}) shows directly that the
    peripheral eigenvalues are closed under powers, and hence are roots of unity
    by Theorem~\ref{thm:peripheral_powers}. It
    requires both normalizations simultaneously. The more general adjoint-fixed-point
    formulation below covers the case where only unitality and a positive definite
    adjoint fixed point are available.
\end{remark}

\subsection{The Kadison--Schwarz inequality and the multiplicative domain}
\label{subsec:ks_input_channels}

\begin{definition}[Kraus maps, adjoints, and normalizations]
    \label{def:kraus_map_channels}
    \lean{KadisonSchwarz.krausMap, KadisonSchwarz.krausAdjointMap,
        KadisonSchwarz.IsUnitalKraus, KadisonSchwarz.IsTPKraus}
    \leanok
    For operators $\{K_i\}_{i=0}^{d-1} \subseteq \MN{D}$, the associated
    \emph{Kraus map} and \emph{adjoint Kraus map} are
    \[
        E(X) = \sum_i K_i X K_i^\dagger,
        \qquad
        E^*(X) = \sum_i K_i^\dagger X K_i.
    \]
    The Kraus family is \emph{unital} when $\sum_i K_i K_i^\dagger = \Id$ and
    \emph{trace-preserving} when $\sum_i K_i^\dagger K_i = \Id$.
\end{definition}

\begin{theorem}[Kadison--Schwarz inequality]
    \label{thm:kadison_schwarz_channels}
    \lean{KadisonSchwarz.kadison_schwarz}
    \leanok
    \uses{def:kraus_map_channels}
    Let $E(X) = \sum_i K_i X K_i^\dagger$ be a unital Kraus map.
    Then for every $X \in \MN{D}$,
    \[
        E(X^\dagger X) \ge E(X)^\dagger E(X).
    \]
    This is \cite[Equation~(5.2)]{Wolf2012Quantum}.
\end{theorem}

\begin{proof}\leanok
    Apply the unital completely positive map entrywise to the positive
    $2 \times 2$ block matrix
    $\bigl(\begin{smallmatrix} X^\dagger X & X^\dagger \\ X & \Id
        \end{smallmatrix}\bigr)$ and take the Schur complement of the lower-right block.
\end{proof}

\begin{theorem}[Hilbert--Schmidt contraction]
    \label{thm:hs_contraction_channels}
    \lean{KadisonSchwarz.hilbertSchmidt_contraction}
    \leanok
    \uses{def:kraus_map_channels, thm:kadison_schwarz_channels}
    If a Kraus map is both unital and trace-preserving, then
    \[
        \tr(E(X)^\dagger E(X)) \le \tr(X^\dagger X)
    \]
    for all $X \in \MN{D}$.
\end{theorem}

\begin{proof}\leanok\uses{thm:kadison_schwarz_channels}
    The Kadison--Schwarz gap is positive semidefinite by
    Theorem~\ref{thm:kadison_schwarz_channels}. Trace preservation then identifies
    the traces of its two terms.
\end{proof}

\begin{theorem}[KS gap decomposition]
    \label{thm:ks_gap_decomp_channels}
    \lean{KadisonSchwarz.ks_gap_eq_sum_squares}
    \leanok
    \uses{def:kraus_map_channels}
    For a unital Kraus map,
    \[
        E(X^\dagger X) - E(X)^\dagger E(X)
        = \sum_i R_i^\dagger R_i,
        \qquad
        R_i := X K_i^\dagger - K_i^\dagger E(X).
    \]
\end{theorem}

\begin{proof}\leanok
    Expand the right-hand side, distribute the products, and use
    $\sum_i K_i K_i^\dagger = \Id$ to recover the Kadison--Schwarz gap.
\end{proof}

\begin{theorem}[KS equality implies Kraus commutation]
    \label{thm:kraus_commute_channels}
    \lean{KadisonSchwarz.kraus_commute_of_ks_equality}
    \leanok
    \uses{thm:ks_gap_decomp_channels}
    Let $E$ be a unital Kraus map.
    If $E(X^\dagger X) = E(X)^\dagger E(X)$, then, for every~$i$,
    \[
        X K_i^\dagger = K_i^\dagger E(X).
    \]
\end{theorem}

\begin{proof}\leanok\uses{thm:ks_gap_decomp_channels}
    The Kadison--Schwarz gap is a sum of positive semidefinite terms
    $\sum_i R_i^\dagger R_i$. If the gap vanishes, each $R_i$ must vanish.
\end{proof}

\begin{theorem}[KS equality for peripheral eigenvectors]
    \label{thm:ks_peripheral_channels}
    \lean{KadisonSchwarz.ks_equality_of_peripheral_eigenvector}
    \leanok
    \uses{def:kraus_map_channels, thm:kadison_schwarz_channels}
    Let $E$ be a Kraus map that is both unital and trace-preserving.
    If $E(X) = \mu X$ with $|\mu| = 1$, then
    \[
        E(X^\dagger X) = E(X)^\dagger E(X).
    \]
\end{theorem}

\begin{proof}\leanok
    The Kadison--Schwarz gap is positive semidefinite by
    Theorem~\ref{thm:kadison_schwarz_channels}. Trace preservation and the relation
    $E(X)=\mu X$ with $|\mu|=1$ show that its trace is zero, hence the gap vanishes.
\end{proof}

\begin{theorem}[Left multiplicative identity from KS equality]
    \label{thm:left_multiplicative_identity_channels}
    \lean{KadisonSchwarz.multiplicative_domain_left}
    \leanok
    \uses{thm:kraus_commute_channels}
    Let $E$ be a unital Kraus map.
    If $E(X^\dagger X) = E(X)^\dagger E(X)$, then for every $Y \in \MN{D}$,
    \[
        E(X^\dagger Y) = E(X)^\dagger E(Y).
    \]
\end{theorem}

\begin{proof}\leanok\uses{thm:kraus_commute_channels}
    The commutation relation of Theorem~\ref{thm:kraus_commute_channels}
    lets one move $X^\dagger$ past each Kraus operator inside the sum.
\end{proof}

\begin{definition}[Multiplicative domain]
    \label{def:multiplicative_domain_channels}
    \lean{KadisonSchwarz.multiplicativeDomain}
    \leanok
    \uses{def:kraus_map_channels}
    The \emph{multiplicative domain} of a Kraus map $E$ is the set of matrices
    that are simultaneously left- and right-multiplicative for $E$.
\end{definition}

\begin{theorem}[Multiplicative domain is a $*$-subalgebra]
    \label{thm:md_star_subalgebra_channels}
    \lean{KadisonSchwarz.multiplicativeDomainStarSubalgebra}
    \leanok
    \uses{def:multiplicative_domain_channels, def:kraus_map_channels}
    For a unital Kraus map, the multiplicative domain is a $*$-subalgebra of
    $\MN{D}$.
\end{theorem}

\begin{proof}\leanok
    The multiplicative identities are stable under addition, multiplication,
    and adjoint, and they contain the identity matrix.
\end{proof}

\subsection{Peripheral closure via adjoint fixed point}%
\label{subsec:peripheral_closure_fixedpoint}

The following results replace the trace-preservation hypothesis by the existence
of a positive definite fixed point for the \emph{adjoint} Kraus map,
matching the weighted-trace argument in
\cite[Appendix~A]{Cirac2016MPDO_arXiv}.

\begin{lemma}[Peripheral closure via adjoint fixed point]%
    \label{lem:peripheral_closed_powers_irred_fp}
    \lean{MPSTensor.peripheralEigenvalues_pow_mem_of_irreducible_unital_of_adjoint_fixedPoint}
    \leanok
    \uses{def:kraus_map_channels, def:irreducible_cp}
    Let $E(X) = \sum_i K_i X K_i^\dagger$ be a Kraus map that is
    unital ($\sum_i K_i K_i^\dagger = \Id$), and assume:
    \begin{enumerate}
        \item the adjoint Kraus map $E^*(X) = \sum_i K_i^\dagger X K_i$
              has a positive definite fixed point $\rho > 0$, and
        \item $E$ is irreducible.
    \end{enumerate}
    Then $\mathrm{peripheral}(E)$ is closed under powers:
    if $\mu \in \mathrm{peripheral}(E)$ then $\mu^n \in \mathrm{peripheral}(E)$
    for every $n \in \N$.
\end{lemma}

\begin{proof}\leanok
    \uses{thm:kadison_schwarz_channels, thm:psd_fp_pd_irred,
        thm:kraus_commute_channels, thm:left_multiplicative_identity_channels}
    Let $E(X) = \mu X$ with $|\mu| = 1$, and set
    $G := E(X^\dagger X) - E(X)^\dagger E(X)$.
    By Kadison--Schwarz, $G \ge 0$.
    Since $E^*(\rho) = \rho$,
    \[
        \tr(\rho G)
        = \tr(\rho E(X^\dagger X)) - \tr(\rho E(X)^\dagger E(X))
        = \tr(E^*(\rho) X^\dagger X) - \tr(\rho E(X)^\dagger E(X))
        = \tr(\rho X^\dagger X) - |\mu|^2 \tr(\rho X^\dagger X)
        = 0.
    \]
    Thus the trace of the Kadison--Schwarz gap against the adjoint
    fixed point vanishes. Because $\rho > 0$ and $G \ge 0$, this forces
    $G = 0$, so $E(X^\dagger X) = E(X)^\dagger E(X)$.
    Hence $X^\dagger X$ is a PSD fixed point, which is
    positive definite by irreducibility
    (Theorem~\ref{thm:psd_fp_pd_irred}).
    So $X$ is invertible.
    By Theorem~\ref{thm:kraus_commute_channels} (KS equality implies Kraus commutation),
    $X$ commutes with each Kraus operator; iterating via the left
    multiplicative identity (Theorem~\ref{thm:left_multiplicative_identity_channels})
    gives $E(X^n) = \mu^n X^n$, so $\mu^n$ is peripheral.
\end{proof}

\begin{theorem}[Roots of unity via adjoint fixed point]%
    \label{thm:peripheral_roots_irred_fp}
    \lean{MPSTensor.peripheral_isRootOfUnity_of_irreducible_unital_of_adjoint_fixedPoint}
    \leanok
    \uses{def:kraus_map_channels, def:irreducible_cp, def:peripheral_eigenvalues}
    Let $E(X) = \sum_i K_i X K_i^\dagger$ be a Kraus map that is
    unital, and assume that the adjoint Kraus map
    $E^*(X) = \sum_i K_i^\dagger X K_i$ has a positive definite fixed
    point and that $E$ is irreducible.
    Then every peripheral eigenvalue of $E$ is a root of unity.
\end{theorem}

\begin{proof}\leanok
    \uses{lem:peripheral_closed_powers_irred_fp, thm:peripheral_powers}
    Combine the closure-under-powers result
    (Lemma~\ref{lem:peripheral_closed_powers_irred_fp})
    with Theorem~\ref{thm:peripheral_powers}.
\end{proof}

\begin{remark}\label{rem:adjoint_fp_more_general}\leanok
    Lemma~\ref{lem:peripheral_closed_powers_irred_fp}
    and Theorem~\ref{thm:peripheral_roots_irred_fp}
    cover the bi-canonical (unital + TP) case as well:
    trace preservation gives
    $E^*(\Id) = \Id$, which is a positive definite fixed
    point of the adjoint.
    The adjoint-fixed-point formulation matches the
    argument in \cite[Appendix~A]{Cirac2016MPDO_arXiv},
    which uses a faithful invariant state rather than
    bi-canonicality.
    Theorem~\ref{thm:peripheral_roots_irred_fp} establishes the
    root-of-unity conclusion in this adjoint-fixed-point formulation.
    In the same adjoint-fixed-point formulation, Chapter~\ref{ch:spectral}
    proves the corresponding cyclic structure: the cyclic group of
    peripheral eigenvalues is Theorem~\ref{thm:peripheral_cyclic_structure},
    and the cyclic projection decomposition is
    Theorem~\ref{thm:cyclic_decomposition_irreducible_schwarz}.
\end{remark}

\begin{definition}[Channel period]\label{def:channel_period}
    \lean{channelPeriod}
    \leanok
    \uses{def:peripheral_eigenvalues}
    The \emph{period} of a channel $E$ is the number of
    peripheral eigenvalues (counted without multiplicity):
    $p(E) := |\mathrm{peripheral}(E)|$.
\end{definition}
